\subsection{What if a current recipient would like to move?}\label{sec:vacancy-chains}
A vacancy can help someone who cannot use it directly. Suppose that Ana and Ben both see only course $x$, which gives Ana priority. Ben goes unplaced. If Ana ranks the vacant course $y$ above $x$, showing her $y$ lets her move there and Ben take $x$. Ben needs no change to his own shortlist. More generally, can we find such trades without making another applicant worse off?

Fix a stable assignment $M$ on the old display $D$, strict applicant preferences and strict, responsive program priorities. For program $j$, its \emph{claimants} are the applicants with $j\in D_i$ who strictly prefer it to $M(i)$. An unmatched applicant prefers every acceptable option to being unmatched. A program with any claimants must already be full and give every incumbent higher priority than every claimant. Write $H_i$ for the hidden options permitted by the intervention. We use either all hidden acceptable options or only those in the applicant's original top five. Keep each baseline contract. Among $D_i\cup H_i$, admit a possible move $(i,j)$ if $i$ strictly prefers $j$ to $M(i)$ and $j$ ranks $i$ at least as high as its highest-priority claimant. A program without claimants imposes no priority restriction. In the quota model, a hidden option may enter if the applicant has a replaceable discretionary slot. Protected options always remain visible.

\begin{theorem}\label{thm:chain-flow}
Let $L$ be the largest number of additional placements in the resulting graph subject to keeping every incumbent matched and each program at least as full as before. A single maximum flow finds $L$ and a collection of vacancy chains, each adding one person. For any subcollection of whole chains, at most one substitution per participating applicant implements a stable witness on the new display. Applicant-proposing DA places exactly the witness's recipients and gives each a weakly preferred assignment. A second maximum flow on the weak improvements in $D_i\cup H_i$, dropping the priority and slot restrictions, gives an upper bound $U$ on additional placements in any assignment weakly improving every applicant on these pairs at the same capacities.
\end{theorem}
\begin{proof}
Start with the flow representing $M$. In its residual network, connect the source to each unmatched applicant, direct available contracts from applicants to programs, reverse each held contract, and connect each program to the sink with its residual capacity. Integral augmentations preserve all current recipients and never reduce program occupancy. The omitted arcs into the source and out of the sink cannot occur on a simple augmenting path. Thus max-flow/min-cut gives the optimum. Discard flow cycles and decompose the remainder into paths. Each applicant lies on at most one path: an unmatched applicant has one unit of inflow from the source, and a matched applicant has one unit from the reverse of their old contract. A program may occur on several paths. Any subset of paths preserves the same inequalities.

Call the resulting assignment $W$. For each hidden witness contract, substitute its program for a displayed discretionary option. An old edge preferred to $W(i)$ was already preferred to $M(i)$, so its applicant was a claimant. That program remains full. Every new holder ranks at least as high as the best old claimant. Equality can occur only if the holder is that claimant, who cannot then block their own contract. Every added edge belongs to $W$, so $W$ is stable. Applicant optimality and invariance of the matched set yield the DA conclusion \citep{gale1962,roth1986}. Finally, every allocation weakly improving on $M$ uses only applicant-approved edges of the second flow graph, or baseline contracts. A maximum ordinary flow there bounds its size; augmenting from $M$ attains that graph's maximum without reducing occupancy or unmatching incumbents.
\end{proof}

\begin{table}[ht]\centering\small
\caption{Additional placements after tentative offers. Direct referrals fix all old contracts. Chains may move incumbents to options they prefer. The ceiling drops program priorities. The right three columns restrict newly displayed options to each person's original top five.}\label{tab:chain-flow}
\begin{tabular}{@{}l rrr rrr@{}}\toprule
 & \multicolumn{3}{c}{All acceptable ranks} & \multicolumn{3}{c}{Top five only}\\
\cmidrule(lr){2-4}\cmidrule(l){5-7}
Starting policy & Direct & Chains & Ceiling & Direct & Chains & Ceiling\\\midrule
NYC ordinary, uniform &215 &641 &1,006 &35 &411 &629\\
NYC ordinary, Haiku &683 &983 &1,353 &204 &610 &843\\
NYC enriched, uniform &518 &1,646 &2,102 &49 &499 &577\\
NYC enriched, Haiku &1,926 &2,493 &2,844 &370 &804 &901\\
Chile, rank-blind &0 &53 &133 &0 &50 &127\\
Chile, rank-informed &409 &475 &649 &1 &2 &4\\\bottomrule
\end{tabular}
\end{table}

In Chile, the rank-informed menus began 1,787 placements behind the rank-blind menus across four cohorts. Even ignoring priorities, giving every incumbent a course at least as preferred as their original assignment leaves a shortfall of at least $1,787-649=1,138$ placements relative to the unchanged rank-blind policy. Requiring newly shown options to be in a pupil's original top five lowers the ceiling from 649 to four. Chile's applicant-proposing replay uses reported courses, pooled capacities, fixed priorities and protected continuity options. The bound concerns this reported option set.

In the rental simulations the corresponding answer changes with the options people accept. With all modeled acceptable pairs available, chains after Haiku reach or exceed unchanged uniform headcount in 226 of 240 ordinary markets and all 240 enriched markets. In the enriched design, requiring newly shown units to be in the renter's top five gives 804 additional placements against an original deficit of 898. Even relaxing the priority constraint yields only 901. When all acceptable options may enter, applying chains to both starting policies leaves Haiku 128 placements behind uniform in ordinary markets and 51 behind in enriched markets. Each design pools 64 simulated renters in each of 240 markets. These are offer counts under modeled rental preferences; neither the identity of the original uniform recipients nor a signed lease follows from recovering the headcount.

Existing mechanisms reassign seats after cancellations \citep{seatReassignment}, study entry in a second round \citep{tworound}, or place additional capacity where it improves assignments \citep{capacityPlanning}. Guaranteed referrals to new schools have been studied in Chile \citep{lepesure}, and allocation facilitators can optimize changes to a compatibility graph \citep{facilitate}. Here the seats, acceptable options and realized priorities stay fixed. The calculation asks how many already acceptable contracts a list of unchanged size can recover under DA, with an explicit bound from applicants' willingness to move. Using a chain requires coordinated take-up by its movers. Selective display can remove an unchosen course on which someone previously had a priority claim. Stability here is assessed on the final displayed lists. The welfare guarantee uses the original ranking of the assigned courses. The theorem gives no incentive guarantee for preference reports.
